\section{Branch A vs.~Branch C --- the tension presented without endorsement}
\label{sec:branch-a-vs-c}

Four open PRs against \texttt{main} make verdicts about $r_e/r_0$ that need
reconciliation, and the reconciliation is exactly what this packet asks the
senior author to provide.

\paragraph{Branch A (Z-universal cutoff) --- PRs \#84, \#85, \#86.}
\begin{itemize}[topsep=2pt,itemsep=0pt]
  \item Reads (III.22)'s closed form $r_e/r_0 = (2 - a_e)/(2(2 + a_e))$ as
        $Z$-independent: $a_e$ is the free-electron anomaly,
        $r_0 = e^2/(m_e c^2)$ depends only on electron mass.
  \item ``The radiating particle is the same electron at every $Z$, so
        $r_e/r_0$ is identical at $Z=1$ and $Z=3$.''
        (\texttt{r\_e\_Li2plus\_joint\_fit.wl}, gating result.)
  \item Verdict: \textbf{Branch A by construction}; Branch~B
        (framework-derivable $Z$-scaling) ruled out structurally.
  \item Empirical context: three of four Li$^{2+}$ comparisons are
        🔴~\textbf{BLOCKED} (mis-provenanced brief measurements ---
        \S\ref{sec:mis-provenance-audit}), so PRs~\#84/\#85/\#86
        cannot compute cross-$Z$ residuals.
\end{itemize}

\paragraph{Branch C (per-$Z$ QED inheritance) --- PR \#87.}
\begin{itemize}[topsep=2pt,itemsep=0pt]
  \item Supplies the cross-$Z$ residual table at $Z = 2, 6, 8, 14, 50$ using
        valid measurements.
  \item Tests (Z-i) universal cutoff: $\chi^2 = 1.0\times 10^{16}$
        $\Rightarrow$ \textbf{Outcome A decisively rejected}.
  \item Per-$Z$ back-fit $x^{(Z)}$ follows
        $x = c_0 + c_2 (Z\alpha)^2 + c_4 (Z\alpha)^4$ with
        $c_0 = x_{\rm univ}$, $c_2 \approx 1/6$ (QED's bound-state coefficient).
  \item Verdict: \textbf{Outcome C} --- $Z$-specific cutoff through
        QED-inheritance. ``The verdict shape is Branch~C, not Branch~A.''
        (PR \#87 body, verbatim.)
\end{itemize}

\paragraph{The non-contradiction reading (substantive AI synthesis).}
The two verdicts can both be honest readings of different things. The
algebra of (III.22), taken literally, has no $Z$ in the coefficients ---
Branch~A is correct as an algebraic statement \emph{about the formula}.
The empirical data across $Z$ rejects fixing a single $x$ value across all
$Z$ --- Branch~C is correct as an empirical statement \emph{about the
fitted cutoff in nature}.

The synthesis: \emph{the framework's apparatus supplies the (III.22) formula
$g_r(x)$ and leaves each state's cutoff free per~(III.23); the apparatus
does not internally derive the binding coefficient $-(Z\alpha)^2/3$. So the
framework, as published, is consistent with Branch~A (its formula has no
$Z$) and the empirical $r_e^{(Z)}/r_0$ inherits QED's bound-state structure
when fit ion-by-ion. The framework's content here is the (III.22)
identification of the cutoff with the anomaly via the $g_r(x)$ formula, not
a derived prediction of $a_e^{\rm bound}(Z\alpha)$.}

\paragraph{Implication for the campaign-closing statement.}
The reading above does not change the honest-scope statement (zero of 28
results constitute independent corroboration distinct from textbook QM); it
sharpens \emph{why}: the (III.22) cutoff--anomaly identification is the
framework's contribution; the empirical $a_e^{\rm bound}(Z\alpha)$ curve is
QED's. Both readings agree that the framework does not internally derive
the binding coefficient. The disagreement is about how to label the
resulting verdict ---
``Branch~A by construction'' (algebraic, Li$^{2+}$-branches) vs.~
``Branch~C QED-inheritance'' (empirical, Z-scan).
The label choice does not change any prediction or honest-scope statement;
it changes what \texttt{10\_CrossComparison.md}~\S1 should say going forward
about the $r_e$ resolution's \emph{character}.

\paragraph{Verdict question for the senior author.}
Both verdicts on \texttt{FINDINGS\_for\_author\_review.md} Finding~2 must
land on \texttt{main} consistently, and the sequencing of merges across the
four BS PRs depends on which label the senior author endorses:
\begin{itemize}[topsep=2pt,itemsep=0pt]
  \item \textbf{If Branch~A by construction is preferred:} PR~\#87's
        $\chi^2 = 10^{16}$ result needs to be re-framed in the merged text
        as ``the per-$Z$ back-fit cutoff curve is QED-inherited; the
        (III.22) algebra produces the universal $r_e/r_0 = x_{\rm univ}$ as
        the $Z\to 0$ limit, and the inheritance curve is empirical content
        outside the framework's algebraic apparatus.''
  \item \textbf{If Branch~C QED-inheritance is preferred:} PRs~\#84/\#85/\#86's
        ``Branch~A by construction'' verdicts need to be qualified as
        ``algebra-correct but not the framework-final verdict; the
        empirical inheritance per~(III.23) is the operative reading.''
\end{itemize}

\authorreview{This is the central reconciliation question of the packet.
Which label is the framework's intended verdict, and how should
\texttt{10\_CrossComparison.md}~\S1 read going forward?}

\paragraph{This packet takes no position.}
Both readings are presented above as honestly as we know how to. The
adjudication belongs to the senior author and is exactly what
\S\ref{sec:questions-for-tepper} asks for.
